% section_blocks.tex — the block decomposition along runs of consecutive quadratic residues (second solver, 19.08)
\subsection*{The blocks: runs of consecutive quadratic residues}\label{sec:blocks}

The two constructions above (the eight-point groups, the seven-point groups) are special cases of a decomposition of $\Pset_{-1}$ that
is exact and, as it turns out, carries the whole strength of the line cover by rows, columns and slope-$(\pm1)$ lines.

For a point $q=(x,y)\in\Pset_{-1}$ put $r_+(q)=x-y$ and $r_-(q)=x+y$ modulo $p$.  Since $r_-^2-r_+^2=4xy\equiv\pm4$ according as
$q\in H(1)$ or $q\in H(-1)$, the two residues of a point are linked; we index everything by
\[
 t(q):=\tfrac14\,r_+(q)^2\in\Fp\qquad\text{for }q\in H(1),\qquad t(q):=\tfrac14\,r_-(q)^2\ \ \text{for }q\in H(-1),
\]
so that $q\in H(1)$ has $\{\tfrac14r_+^2,\tfrac14r_-^2\}=\{t,t+1\}$ and $q\in H(-1)$ has $\{\tfrac14r_-^2,\tfrac14r_+^2\}=\{t,t+1\}$:
every point determines an unordered pair of \emph{consecutive} values $t,t+1$, both of which are squares (as $4t=r^2$).  Call
$\{t,t+1\}$ the \emph{edge} of $q$.

\begin{lemma}[edges]\label{lem:edges}
Let $t\in\Fp$ be such that $t$ and $t+1$ are both squares (with $0$ counted as a square).  The points of $\Pset_{-1}$ with edge
$\{t,t+1\}$ are the four lifts of each of the $H(1)$-classes $\kappa$ with $d_\kappa^2\equiv4t$ together with the four lifts of each of
their images under $R$, which are $H(-1)$-classes.  For $t\ne0,-1$ there are four such $H(1)$-classes (two $\sigma$-pairs, for
$d_\kappa\equiv\pm2\sqrt t$), hence $32$ points; for $t=0$ there are two (the $\tau$-fixed classes $(1,1)$, $(-1,-1)$) and for $t=-1$ two
(the $\sigma$-fixed classes, present iff $-1$ is a square), hence $16$ points.  If one of $t,t+1$ is a non-square there are no such points.
The reflections $R$ and $R'$ preserve the edge, so every row and every column of the box consists of four points with the same edge.
\end{lemma}
\begin{proof}
An $H(1)$-class $\kappa=(a,1/a)$ has $r_+\equiv a-1/a=:d_\kappa$ and $r_-\equiv a+1/a=:e_\kappa$ with $e_\kappa^2-d_\kappa^2=4$, so its edge is
$\{\tfrac14d_\kappa^2,\tfrac14d_\kappa^2+1\}$; conversely $d_\kappa\equiv\pm2\sqrt t$ determines $a$ up to the two roots of $a^2-d_\kappa a-1$
(which exist because $d_\kappa^2+4=4(t+1)$ is a square), and those two roots form a $\sigma$-pair, so that the two signs of $d_\kappa$ give
four classes unless $d_\kappa\equiv-d_\kappa$, i.e.\ $t=0$, when $a^2\equiv1$ gives the single pair $\{(1,1),(-1,-1)\}$; the case $t=-1$ is
$e_\kappa\equiv0$, i.e.\ $a^2\equiv-1$.  By $R(x,y)=(p-x,y)$ we have $(r_+,r_-)\mapsto(-r_-,-r_+)$ and by $R'(x,y)=(x,2p-y)$ we have
$(r_+,r_-)\mapsto(r_-,r_+)$: both preserve the pair $\{\tfrac14r_+^2,\tfrac14r_-^2\}$ and exchange $H(1)$ with $H(-1)$, so the classes are as
stated and rows ($\kappa\cup R\kappa$) and columns ($\kappa\cup R'\kappa$) stay inside the edge (Lemma~\ref{lem:rowsRR}).
\end{proof}

Let $Q:=\{t\in\Fp:\ t\ \text{is a square}\}$ (including $t=0$) and consider the graph on $Q$ whose edges are the pairs $\{t,t+1\}$ of
consecutive squares (Lemma~\ref{lem:edges}).  Its connected components are the maximal \emph{runs} of consecutive
squares $t_0,t_0+1,\dots,t_0+k-1$ (with $0$ joining the runs that end at $-1$ and start at $1$, if $-1$ is a square).  For a run $\varrho$
let $B_\varrho\subseteq\Pset_{-1}$ be the union of the points of its $k-1$ edges, so that $|B_\varrho|=32(k-1)$ unless the run contains
$0$, when each of its at most two degenerate edges contributes $16$ instead of $32$.

\begin{theorem}[block decomposition]\label{thm:blocks}
The sets $B_\varrho$, $\varrho$ a run of consecutive squares, partition $\Pset_{-1}$; each $B_\varrho$ is a union of rows and of columns;
and every row, every column and every slope-$(\pm1)$ line of the box lies inside a single $B_\varrho$.  Consequently
\begin{equation}\label{eq:blocks}
 \alpha(\Pset_{-1})\ \le\ \sum_{\varrho}\beta(B_\varrho),
\end{equation}
where $\beta(B)$ is the largest size of an $S\subseteq B$ with at most two points on every row, column and slope-$(\pm1)$ line,
and the same sum with $\beta$ replaced by the value of the linear relaxation equals the value of the linear relaxation for
$\Pset_{-1}$ with rows, columns and slope-$(\pm1)$ lines.  (Lines of other slopes may carry three points of $\Pset_{-1}$ --- two from one
hyperbola and one from the other --- and such a line need not stay inside a block; dropping those constraints is what makes the bound
block-wise, and it only weakens it.)
\end{theorem}
\begin{proof}
Every point lies on exactly one edge (Lemma~\ref{lem:edges}), and the edges of one run are pairwise disjoint sets of classes, so the
$B_\varrho$ partition $\Pset_{-1}$; rows and columns stay inside an edge, hence inside a block.  A line of slope $+1$ with residue $d$
carries the $H(1)$-classes with $d_\kappa\equiv\pm d$ --- edge $\{t,t+1\}$, $t=\tfrac14d^2$ --- and the $H(-1)$-classes with $r_+\equiv\pm d$,
whose edge is $\{t-1,t\}$ (their $r_-^2=d^2-4$); both edges contain the vertex $t$, so they lie in the same run and the line lies in one
block.  Slope $-1$ is symmetric under $R$.  Since the blocks partition $\Pset_{-1}$, a lawful $S$ satisfies
$|S|=\sum_\varrho|S\cap B_\varrho|$, and each $S\cap B_\varrho$ has at most two points on every row, column and slope-$(\pm1)$ line, so
$|S\cap B_\varrho|\le\beta(B_\varrho)$; the constraint matrix of the linear relaxation over those three families is block diagonal.
\end{proof}

% lemma_run.tex — the local law along a run of quadratic residues and the descent formula for signature densities
% (first solver, 2026-08-19).  For the section on the neighbour-graph components (k=-1).  Uses \Fp, amsthm; labels lem:run, prop:descent.

\begin{lemma}[equidistribution along a run]\label{lem:run}
Fix $k\ge2$ and let $C_k$ be the affine curve
$\{(u_1,\dots,u_k):\ u_{i+1}^2=u_i^2+1\ (1\le i<k)\}$, i.e.\ the multiquadratic cover
$\Fp(t)\bigl(\sqrt t,\sqrt{t+1},\dots,\sqrt{t+k-1}\bigr)$ of the $t$-line ($t=u_1^2$).  It is absolutely irreducible of genus
$1+2^{k-2}(k-3)$, and for every $\mathbf h=(h_1,\dots,h_k)\in\Fp^k\setminus\{0\}$ the function $\sum_ih_iu_i$ is non-constant on $C_k$.
Consequently, for all intervals $I_1,\dots,I_k\subseteq(-p/2,p/2)$ and all $\varepsilon_0,\varepsilon_k\in\{\pm1\}$,
\[
 \#\Bigl\{(u_i)\in C_k(\Fp):\ X(u_i)\in I_i\ \forall i,\ \chi(u_1^2-1)=\varepsilon_0,\ \chi(u_k^2+1)=\varepsilon_k\Bigr\}
 =\frac p4\prod_{i=1}^k\frac{|I_i|}p+O_k\bigl(\sqrt p\log^kp\bigr),
\]
where $\chi$ is the quadratic character and $X$ the centred representative.  In particular the positions $X(u_i)/p$ are asymptotically
independent and uniform on $(-\frac12,\frac12)$, and independent of the two Legendre conditions that make the run maximal.
\end{lemma}
\begin{proof}
Irreducibility: the $k$ quadratic extensions are independent because the polynomials $t,t+1,\dots,t+k-1$ are pairwise coprime and squarefree
(for $p>k$), so the Galois group over $\overline\Fp(t)$ is $(\Z/2)^k$ and the cover is connected.  Genus: the cover of degree $2^k$ is
ramified exactly over the $k$ finite points $t=0,-1,\dots,-(k-1)$ and over $\infty$, with inertia $\Z/2$ at each (at $\infty$ the inertia is
generated by the simultaneous flip of all square roots, each $\sqrt{t+i}$ having odd polar order); each such point has $2^{k-1}$ preimages
with $e=2$, so $2g-2=2^k(-2)+(k+1)2^{k-1}$, i.e.\ $g=1+2^{k-2}(k-3)$.  Non-constancy: the group $(\Z/2)^k$ acts on $C_k$ by
$u_i\mapsto\pm u_i$ independently; if $\sum_ih_iu_i$ were constant, averaging over the subgroup flipping only $u_j$ would give $h_ju_j$
constant, hence $h_j=0$ for every $j$.  Now expand the interval indicators in Fourier series (or Selberg polynomials) of $\ell^1$-norm
$O(\log p)$ each and the two Legendre conditions as $\frac12(1\pm\chi)$; the main term is the product of the volumes times
$|C_k(\Fp)|/4=p/4+O_k(\sqrt p)$, and every other term is a mixed character sum
$\sum_{C_k(\Fp)}\chi\bigl((u_1^2-1)^{a}(u_k^2+1)^{b}\bigr)e\bigl(\sum_ih_iu_i/p\bigr)$ with $(a,b,\mathbf h)\ne0$, which is $O_k(\sqrt p)$ by the
Bombieri--Perel'muter bound: the additive part is non-constant when $\mathbf h\ne0$, and for $\mathbf h=0$ the Kummer sheaf is geometrically
non-trivial because $u_1^2-1=t-1$ and $u_k^2+1=t+k$ have odd order at their zeros on the $t$-line and their zeros are unramified in $C_k$.
\end{proof}

\begin{proposition}[signature densities are descent statistics]\label{prop:descent}
Let a component of the neighbour graph be a maximal run $t,t+1,\dots,t+k-1$ of quadratic residues, and for $1\le i\le k-1$ let
$\xi_i\in\{0,1\}$ record whether the $H(1)$-pair carried by the edge $(t+i-1,t+i)$ shares its centre.  Then, as $p\to\infty$,
\[
 \Pr\bigl[(\xi_1,\dots,\xi_{k-1})=\varepsilon\bigr]=\frac{D_k(S_\varepsilon)}{k!}+O_k\bigl(p^{-1/2}\log^kp\bigr),\qquad
 S_\varepsilon=\{i:\varepsilon_i=1\},
\]
where $D_k(S)$ is the number of permutations of $\{1,\dots,k\}$ with descent set exactly $S$.  The types of the vertices are determined by
$\varepsilon$: the vertex between the edges $i-1$ and $i$ is a good group $(4,8,4)$ iff $(\xi_{i-1},\xi_i)=(0,1)$, a $(2,6,6,2)$-group iff
$(1,0)$, and a seven-group otherwise; the end vertices are $(4,2,2)$ iff $\xi_1=1$ (left), resp.\ $\xi_{k-1}=0$ (right), and $(3,3,1,1)$ otherwise.
\end{proposition}
\begin{proof}
Write $u_i$ for a square root of $t+i-1$.  The $H(1)$-pair of the edge $(t+i-1,t+i)$ is $\{a,1/a\}$ with
$a=u_i+u_{i+1}$, $1/a=u_{i+1}-u_i$ (indeed $a\cdot(1/a)=u_{i+1}^2-u_i^2=1$ and $a-1/a=2u_i$ is the residue $d$ of the group), and the pair
shares its centre iff $X(a)X(1/a)<0$; this is independent of the choice of signs of the $u_i$'s.  Put $v_i=X(u_i)/p$.  Since the sign of the
centred representative of $x$ equals the sign of $\sin2\pi x$ and
$\sin2\pi(v+w)\,\sin2\pi(w-v)=\tfrac12\bigl(\cos4\pi v-\cos4\pi w\bigr)$, we get the exact identity
\[
 \xi_i=1\iff\cos4\pi v_i<\cos4\pi v_{i+1}\iff g(v_i)>g(v_{i+1}),\qquad g(x):=\min\bigl(|x|,\tfrac12-|x|\bigr).
\]
By Lemma~\ref{lem:run} the $v_i$ are asymptotically independent and uniform, hence the $g(v_i)$ are independent and uniform on $(0,\frac14)$
(the map $g$ is four-to-one and piecewise linear), so $(\xi_i)$ is the descent pattern of $k$ i.i.d.\ uniform variables, whose law is
$D_k(S)/k!$.  The type dictionary is Lemma~\ref{lem:partners} together with the fact that the pair on an edge is seen as an
$H(1)$-pair by its left vertex and as (the $R$-image of) a $\tau$-pair by its right vertex, which exchanges $A\cup D$ with $B\cup C$.
\end{proof}


The saving $s(\varepsilon):=|B_\varrho|/2-\beta(B_\varrho)$ of a block depends only on its signature (equivalently, on $\varepsilon$): this was
verified for every block occurring in the primes $p\le20000$.  Combining Proposition~\ref{prop:descent} with the exact values of
$\beta$ (a finite computation: \texttt{slack/t221/fast\_blocks.py} constructs the block of a run directly from the square roots
$u_i=\sqrt{t+i-1}$ and \texttt{scipy} solves the integer program), and using that the number of runs of length $k$ is
$2^{-(k+2)}p+O(\sqrt p\log p)$ (Weil), we obtain the following.

\begin{corollary}\label{cor:blockconst}
For every $K\ge2$,
\[
 \alpha(\Pset_{-1})\ \le\ \Bigl(4-\sum_{k=2}^{K}2^{-(k+2)}\sum_{\varepsilon\in\{0,1\}^{k-1}}\frac{D_k(S_\varepsilon)}{k!}\,s(\varepsilon)+o(1)\Bigr)(p-1),
\]
a finite computation for each $K$.  The table of $s(\varepsilon)$ was computed in three independent ways: by constructing the block of every run occurring at the primes
$p\le20000$ (\texttt{slack/t221/fast\_blocks.py}), by an independent implementation through the $\xi$-bits at $p\le9000$
(\texttt{slack/sig\_savings.py}), and --- with no prime at all --- from the \emph{standard model} of a block, in which the positions
$v_i=X(u_i)/p$ are replaced by generic rationals with the prescribed descent pattern and the incidences are decided by exact arithmetic
(\texttt{slack/t221/standard\_model.py}); the three agree wherever they overlap ($500$ comparisons for $k\le9$, no discrepancy), and the
standard model supplies every $\varepsilon$, including those too rare to occur at moderate primes.  With $K=9$ the sum equals
$68350729/123863040$, so
\[
 \alpha(\Pset_{-1})\ \le\ \frac{427101431}{123863040}\,(p-1)+o(p)\ =\ (3.44817\ldots+o(1))\,(p-1).
\]
Since $s(\varepsilon)\le4(k-1)$, the terms with $k>9$ contribute at most $\sum_{k>9}2^{-(k+2)}4(k-1)=0.0196$, so the constant of this
certificate lies in $[3.4286,\,3.4482]$; numerically it is $3.4413\ldots$
\end{corollary}

\begin{remark}\label{rem:blocks}
(a) The blocks are small: the run through $t$ has length $k$ with probability $\asymp2^{-k}$, so $\beta(B_\varrho)$ is a bounded computation,
and by Weil's bound the number of runs of a given length and given \emph{type} (the sequence of line-size profiles of the groups at the
interior vertices, which are the $(4,8,4)$, $(3,7,5,1)$, $(2,6,6,2)$ groups of Section~\ref{sec:two}, the end vertices carrying two-class
groups) is $c\,p+O(\sqrt p\log^Cp)$ with an explicitly computable $c$ --- the local law of Propositions~\ref{prop:m8asym},
\ref{prop:m7asym} applied along the run.  Theorem~\ref{thm:two} is the case of a single $(4,8,4)$ interior vertex ($\beta=|B|/2-8$ for the
two edges at it), Theorem~\ref{thm:seven} the case of a $(3,7,5,1)$ vertex ($\beta=|B|/2-4$).
(b) Numerically (\texttt{slack/t221/}, all primes $200\le p\le1400$): the saving $|B_\varrho|/2-\beta(B_\varrho)$ depends only on the type of the run
(no exception among the several hundred blocks tested), and $\sum_\varrho(|B_\varrho|/2-\beta(B_\varrho))=0.5587\,p$ on average over
$401\le p\le1103$, in agreement with Corollary~\ref{cor:blockconst}; truncating to runs of length $\le L_0$ gives
$3.595$ for $L_0=5$, $3.480$ for $L_0=7$ and $3.452$ for $L_0=10$.  At a given prime the bound~\eqref{eq:blocks} is at most the value of the linear relaxation with
all these lines, and strictly below it when integrality bites inside a block ($3.4286$ vs $3.4286$ at $p=113$, $3.3824$ vs $3.3873$ at
$p=137$, $3.4343$ vs $3.4545$ at $p=199$, $3.4100$ vs $3.4300$ at $p=401$); the \emph{proved} constant $3.4616$ of
Corollary~\ref{cor:blockconst} is larger than the numerical value $3.4413$ of the same certificate only because of the truncation at $K=8$.
(c) The method is special to $k=\pm1$.  For $\Pset_k=H(1)\cup H(k)$ a point of $H(c)$ has $r_-^2-r_+^2=4c$, so with $t=\tfrac14r_+^2$ the points of
$H(1)$ join $t$ to $t+1$ and those of $H(k)$ join $t$ to $t+k$: for $k=\pm1$ the two rules coincide and the neighbour graph on the squares has mean
degree $1$ --- subcritical, with components of size $O(\log p)$ --- whereas for every other $k$ the mean degree is $2$, the critical value, and the
components have size about $p^{2/3}$ (at $p=4001$: largest component $11$ for $k=-1$ against $22$, $32$, $95$, $177$ for $k=2,3,5,7$;
\texttt{slack/blocks\_general\_k.py}).  So a finite table of block optima exists only for the pair $H(1)\cup H(-1)$; for a general second hyperbola
the decomposition is still valid but useless, and Section~\ref{sec:general} (a saving of order $\sqrt p$) is what remains.
(d) The decomposition explains the plateau of Section~\ref{sec:beyond}: no certificate built from rows, columns and slope-$(\pm1)$ lines can
give less than $\sum_\varrho\beta(B_\varrho)$, and this is $\approx3.44(p-1)$, far above the computed maxima $3(p-1)+O(1)$.  The remaining
gap is carried by the three-point lines of the other slopes, which are invisible to the block decomposition.
\end{remark}
